\documentclass[twocolumn,a4paper,10pt]{ctexart}
\usepackage[top=2.2cm,bottom=2.2cm,left=1.8cm,right=1.8cm,columnsep=0.7cm]{geometry}
\setlength{\parindent}{2em}
\usepackage{amsmath,amssymb,bm,mathtools}
\usepackage{graphicx,booktabs,caption,float}
\usepackage[colorlinks,linkcolor=black,citecolor=black,urlcolor=blue]{hyperref}
\usepackage{listings,xcolor}
\usepackage{fancyhdr}
\pagestyle{fancy}\fancyhf{}
\fancyhead[L]{\small 系统工程与电子技术}
\fancyhead[R]{\small Systems Engineering and Electronics}
\fancyfoot[C]{\thepage}
\renewcommand{\headrulewidth}{0.4pt}
\graphicspath{{/mnt/linuxdata/cuda_ik_accelerateion/standard_robot_cuda_ik/experiments/figures/}}
\lstset{
  basicstyle=\ttfamily\footnotesize,
  commentstyle=\color{green!45!black}\itshape,
  keywordstyle=\color{blue}\bfseries,
  stringstyle=\color{orange},
  numbers=left,
  numberstyle=\tiny\color{gray},
  stepnumber=1,
  numbersep=5pt,
  backgroundcolor=\color{gray!5},
  frame=single,
  framerule=0.5pt,
  rulecolor=\color{gray!50},
  breaklines=true,
  breakatwhitespace=false,
  tabsize=3,
  language=C++
}
\DeclareMathOperator{\SE}{SE}\DeclareMathOperator{\SO}{SO}
\newcommand{\ud}{\mathrm{d}}\newcommand{\R}{\mathbb{R}}
\begin{document}

\twocolumn[%
{\centering
{\LARGE\bfseries 基于 CUDA 小矩阵加速的机械臂批量逆运动学求解}\par\vspace{3mm}
{\large 匿名稿（通讯作者信息在单独标题页）}\par\vspace{2mm}
{\small （作者单位信息在单独标题页）}\par\vspace{3mm}
{\small\bfseries 文章编号：}（预留）\quad{\small\bfseries DOI：}（预留）\par\vspace{2mm}
\begin{minipage}{0.9\textwidth}\small
{\bfseries 摘　要：} 针对机械臂批量逆运动学求解中 $6\times6$ 小矩阵反复求解与多 kernel 调度带来的固定开销问题，提出基于计算统一设备架构（compute unified device architecture，CUDA）底层优化的批量阻尼最小二乘法加速方法。将每个目标映射为 1 个 CUDA 线程块，在单次 kernel launch 中完成全部迭代；设计寄存器驻留的 $6\times6$ $\mathrm{LDL}^\top$ 求解器替代 cuBLAS；采用 FP32/FP64 混合精度策略。以官方 UR10 为对象构建统一 benchmark，从端到端、GPU stream 和有效吞吐三个维度报告性能并引用 cuRobo（CUDA Graph 开启）为公平基线。实验表明：当 $N\ge500$ 时本文方法吞吐稳定在 128k--157k targets/s，GPU 时间线性扩展（$R^2>0.999$）；在 $N=100$ 极小批量下 GPU 流延迟较 cuRobo-Graph 降低约 52\%。当 $N\ge500$ 时 cuRobo-Graph 依靠 Graph 和批处理优化通常获得更高吞吐，但其存在批量依赖的非单调退化。本文方法的优势在于 kernel launch 数恒为 1、结构简单、批量扩展可预测。Jacobian 精度分析表明 FP32 混合精度在 $\varepsilon=10^{-6}$ 下对应约 4--5 mm 的 IK 收敛精度边界——适用于中等精度批量采样与候选解筛选。\par\vspace{2mm}
{\bfseries 关键词：} 逆运动学；计算统一设备架构；小矩阵求解；混合精度；Bank 冲突\par
{\bfseries 中图分类号：} TP391.4\quad{\bfseries 文献标志码：} A\par
{\bfseries 收稿日期：}（预留）；{\bfseries 修回日期：}（预留）；{\bfseries 网络优先出版日期：}（预留）。\par
\end{minipage}\par
}]

{\bfseries\large CUDA-Accelerated Small-Matrix Solver for Batch Inverse Kinematics of Robot Manipulators}

\vspace{1mm}
{\small （Author information on separate title page）}

\vspace{2mm}
{\bfseries Abstract:} This paper presents a CUDA-based low-level acceleration method for batch damped least-squares inverse kinematics (IK). Each target is mapped to one CUDA block (128 threads), with all iterations in a single kernel launch. A register-resident $6\times6$ $\mathrm{LDL}^\top$ solver replaces cuBLAS. A mixed-precision strategy applies FP32 to forward kinematics and Jacobian evaluation while retaining FP64 for $\mathrm{LDL}^\top$ and convergence checks. On a unified UR10 benchmark with cuRobo (CUDA Graph enabled) as the fair baseline, the proposed method achieves 128k--157k targets/s for $N\ge500$ with strictly linear GPU time scaling ($R^2>0.999$). At $N=100$, GPU stream latency is $\approx$52\% lower than cuRobo-Graph. For $N\ge500$, cuRobo-Graph generally achieves higher throughput via Graph and internal batching, but exhibits batch-size-dependent non-monotonic degradation. The advantage of the proposed method lies in its structural simplicity: constant-1 kernel launch count and predictable batch scaling. Jacobian precision analysis reveals an IK convergence precision boundary of $\approx$4--5 mm under FP32 mixed precision, suitable for medium-precision batch sampling and candidate screening.

\vspace{1mm}
{\bfseries Keywords:} inverse kinematics; CUDA; small-matrix solver; mixed precision; bank conflict

\section{引 言}

机械臂批量逆运动学（inverse kinematics，IK）求解是轨迹优化、料箱抓取和采样规划等应用的核心计算环节[1]。给定 $N$ 个末端目标位姿，批量 IK 需为每个目标求解关节角使正运动学（forward kinematics，FK）输出匹配目标。尽管部分具有球腕结构的 6R 工业机械臂可推导解析 IK，数值 IK 因模型适应性强、易处理关节限位和扩展约束，在批量采样与通用软件框架中广泛采用[2-3]。贾龙飞等[4]综述了三大类 IK 求解方法，指出数值解法通用性好但计算量大。批量 IK 的计算特征：（1）每次迭代须反复求解 $6\times6$ 极小固定规模线性系统；（2）每目标需 10--80 次迭代收敛；（3）目标数 $N$ 可达数千至上万。

{\bfseries 批量 IK 计算方法。} Orocos KDL 提供成熟的 C++ FK/IK 工具链[5]，但串行迭代在批量下吞吐有限。Buss 等[2]系统比较了 Jacobian 转置、伪逆和阻尼最小二乘法（damped least squares，DLS），指出 DLS 在奇异点附近数值稳定性更优。

{\bfseries GPU 加速逆运动学。} NVIDIA cuRobo[6]将 IK 建模为非线性优化，在 GPU 上并行评估粒子并通过 L-BFGS 求解，是当前最具影响力的 GPU IK 框架。HJCD-IK[7]结合贪婪坐标下降与 Jacobian 伪逆精化。ManipulaPy[8]基于指数积公式和 CuPy 自定义 CUDA kernel。Kino-PAX[9]等 GPU 运动规划工作的``全 GPU 原生''架构理念启发了本文的单 kernel 设计。上述方法大多基于 PyTorch、JAX 或 CuPy 等高层框架，其多 kernel 调度和跨语言调用栈在小粒度批量 IK 中引入了不可忽略的固定开销。

{\bfseries CUDA 底层优化。} 面向极小矩阵（$n\le32$）的 CUDA 优化在数值线性代数领域已有深入研究。Abdelfattah 等[10]针对百万级批量的 LU/QR/Cholesky 分解，通过寄存器分块实现相对厂商库 11.8 倍加速。刘世芳等[11]在 GPU 上实现批量 LU 矩阵求逆，达到 cuBLAS 的 13 倍加速。Shridhar[12]在 RTX 3080 上证实自定义小矩阵求逆 kernel 仅占总运行时间 10\%，而 cuBLAS 调用占 67\%。Abdelfattah 等[13]利用 Tensor Core 进行三层迭代精化，在关键路径保留 FP64 的前提下实现 3--5 倍加速。

本文核心观点：对 $6\times6$ 固定规模 IK 小矩阵，{\bfseries 手写 CUDA 底层实现可通过减少通用库调度开销和框架层间通信，在极小批量下获得更低延迟，并在全量程保持稳定线性扩展}。贡献：（1）1 block/target 映射与单 kernel 封装，kernel launch 数恒为 1；（2）寄存器驻留 $6\times6$ $\mathrm{LDL}^\top$ 求解器；（3）PaddedMat$6\times8$ 布局将 Bank 冲突降低 63\%；（4）FP32/FP64 混合精度策略，通过 Jacobian 精度验证量化其适用范围边界。

\section{批量 IK 问题建模与 DLS 方法}\label{sec:model}

\subsection{运动学模型}

UR10 为 6 自由度全旋转关节串联机械臂，DH 参数遵循 Universal Robots 官方 URDF[14]。正运动学将关节角 $\mathbf{q}=[q_1,\ldots,q_6]^\top$ 映射为末端齐次变换矩阵 $\mathbf{T}_{\mathrm{ee}}(\mathbf{q})\in\mathrm{SE}(3)$：
\begin{equation}
\mathbf{T}_{\mathrm{ee}}(\mathbf{q}) = \prod_{i=1}^{6} \mathbf{T}_{i-1,i}(q_i)\label{eq:fk}
\end{equation}
实用中采用 Rodrigues 公式计算旋转矩阵。

对目标位姿 $\mathbf{T}_{\mathrm{tgt}}$，定义位姿误差 $\mathbf{e}(\mathbf{q})\in\mathbb{R}^6$，引入对角权重矩阵 $\mathbf{W}$。

\subsection{DLS 迭代与批量特征}

DLS[2-3]在第 $k$ 次迭代中极小化正则化目标：
\begin{equation}
\min_{\Delta\mathbf{q}} \frac{1}{2}\|\mathbf{W}\mathbf{J}\Delta\mathbf{q} + \mathbf{W}\mathbf{e}\|^2 + \frac{1}{2}\lambda\|\Delta\mathbf{q}\|^2\label{eq:dls}
\end{equation}
雅可比由 FK 中心差分计算：
\begin{equation}
\mathbf{J}_{:,j} = \frac{\mathbf{x}(\mathbf{q}+\varepsilon\mathbf{e}_j)-\mathbf{x}(\mathbf{q}-\varepsilon\mathbf{e}_j)}{2\varepsilon},\quad \varepsilon=10^{-6}\ \mathrm{rad}\label{eq:jacobian}
\end{equation}
令式(\ref{eq:dls})梯度为零得阻尼正规方程，记 $\mathbf{H}=\mathbf{J}^\top\mathbf{W}^2\mathbf{J}+\lambda\mathbf{I}$，$\mathbf{g}=\mathbf{J}^\top\mathbf{W}^2\mathbf{e}$，解 $\Delta\mathbf{q}=-\mathbf{H}^{-1}\mathbf{g}$。收敛条件：$\|\mathbf{e}_p\|<10$ mm 且 $\|\mathbf{e}_R\|<5^\circ$，$K_{\max}=160$。

批量计算特征：$N$ 个目标共享模型参数但各自独立状态向量——完全数据并行且计算粒度极小（单目标每迭代约 1 500 FLOP）。

\section{CUDA 小矩阵加速设计}\label{sec:design}

\subsection{1 Block/Target 并行映射与单 Kernel 封装}

采用 $\mathrm{Grid}(N,1,1)+\mathrm{Block}(128,1,1)$ 映射，如图~\ref{fig:arch}。选择 128 线程/block：寄存器用量 94--98/线程，每 SM 可驻留约 5 block；$\mathbf{H}$ 矩阵 36 元素决定最大并行度，128 线程充足。

\begin{figure}[ht]
\centering
\includegraphics[width=\columnwidth]{fig1_architecture.png}
\caption{1 block/target CUDA 并行映射与单 kernel 全迭代结构}
\label{fig:arch}
\end{figure}

DLS 迭代线程分工见表~\ref{tab:thread}。

\begin{table}[ht]
\centering
\caption{DLS 迭代各阶段线程分工}
\label{tab:thread}
\begin{tabular}{lcc}
\toprule
计算阶段 & 线程范围 & 并行度 \\
\midrule
FK/位姿误差/收敛判定 & $t=0$ & 1 \\
数值 Jacobian & $0\le t<6$ & 6 \\
自适应阻尼更新 & $t=0$ & 1 \\
$\mathbf{H}$ 矩阵构造 & $0\le t<36$ & 36 \\
$\mathbf{g}$ 向量构造 & $0\le t<6$ & 6 \\
$\mathrm{LDL}^\top$ 求解/步长钳位 & $t=0$ & 1 \\
关节更新+限位投影 & $0\le t<6$ & 6 \\
\bottomrule
\end{tabular}
\end{table}

{\bfseries 单 Kernel 全迭代封装——与 cuRobo 结构性对比。} 全部 $K_{\max}=160$ 次迭代在单次 \texttt{ik\_batch\_solve<<<N,\;128>>>} 中完成，CPU 端仅需一次同步。cuRobo 执行路径则经历 Python 调用$\to$tensor 分配$\to$子批次划分$\to$多 kernel 启动$\to$L-BFGS 迭代$\to$结果回传。Nsight Systems 实测 $N=4\,000$ 时 cuRobo-NoGraph 产生 14 108 次 kernel launch，启动开销累积 70--140 ms。本文方案 kernel launch 数恒为 1——所有 block 在同一 launch 中并行发射。需指出：总计算量仍随 $N$ 和 $K$ 线性增长，但固定调度开销被压缩至常数级。

精简伪代码：

\begin{lstlisting}[language=C++, caption={DLS单次迭代伪代码（per-block，128线程）}]
// s_q[6], s_T_tgt[16], s_lambda in shared memory
if (threadIdx.x == 0)               // FK computation
    forward_kinematics(s_q, s_T_ee);
__syncthreads();
if (threadIdx.x == 0)               // convergence check
    converged = (pos_err<0.01) && (rot_err<0.0873);
__syncthreads();
if (threadIdx.x < 6)                // numerical Jacobian (6 cols)
    J_col = (FK(q+eps)-FK(q-eps))/(2*eps);
__syncthreads();
if (threadIdx.x == 0)               // adaptive damping
    update_lambda(s_lambda, pos_err);
__syncthreads();
if (threadIdx.x < 36)               // H = J^T W^2 J + lambda*I
    H[row][col] = inner_product(J,W2,row,col);
__syncthreads();
if (threadIdx.x < 6)                // g = J^T W^2 e
    g[tid] = dot_product(J,W2,e);
__syncthreads();
if (threadIdx.x == 0)               // register LDLT + step clamp
    ldlt_solve_6x6(H, g, dq);
__syncthreads();
if (threadIdx.x < 6)                // joint update + limits
    s_q[tid] = clamp(s_q[tid]+dq[tid], q_min[tid], q_max[tid]);
\end{lstlisting}

\subsection{寄存器级 $6\times6$ $\mathrm{LDL}^\top$ 求解器}

cuBLAS 批量接口针对 $n\ge32$ 优化，对 $6\times6$ 矩阵调用开销（$\sim$10--15 $\mu$s）是求解本身（$\sim$0.1 $\mu$s）的 100 倍以上[12]。$\mathrm{LDL}^\top$（$\mathbf{H}=\mathbf{L}\mathbf{D}\mathbf{L}^\top$）避免 Cholesky 所需的 6 次 sqrt（GPU 上 sqrt 吞吐为 FMA 的 1/4--1/8[15]）。对 $n=6$，分解和三角回代编译期完全展开为寄存器级直线代码，算术规模百级标量操作。Nsight Compute 实测：94--98 regs/thread，零 local memory spill。

\subsection{PaddedMat$6\times8$ 共享内存 Bank 冲突降低}

GPU 共享内存以 32 Bank 组织，FP64 元素跨 2 Bank。stride=6 下 $\gcd(12,32)=4$ 导致 2--3 路冲突。stride$6\to8$：行步长 64 字节=16 Bank，$\gcd(16,32)=16$，偶数行 Bank 0--15 与奇数行 Bank 16--31 完全分离。NCU 实测 Bank 冲突从 3 522 降至 1 295 次（$-63\%$，见后文图~\ref{fig:ncu}）。

\subsection{FP32/FP64 混合精度计算路径与精度边界}

Ada Lovelace 的 FP64:FP32 吞吐比 1:64[15]。本文选 FP32+FP64 方案：FP32 提供 6--7 位十进制精度用于 FK/Jacobian，FP64 保留 $\approx10^{-15}$ 精度裕度用于 $\mathrm{LDL}^\top$ 和收敛判定。

\begin{table*}[ht]
\centering\small
\caption{混合精度分配策略}
\label{tab:mixedprec}
\begin{tabular}{lcl}
\toprule
计算阶段 & 精度 & 理由 \\
\midrule
FK 链式乘法 & FP32 & 吞吐为 FP64 的 64 倍 \\
Jacobian 各列 & FP32 & $\varepsilon=10^{-6}$ 下中位误差 $4.2\times10^{-2}$ \\
$\mathbf{H}$ 累加 & {\bfseries FP64} & 36 个内积抑制截断误差 \\
$\mathbf{g}$ 累加 & {\bfseries FP64} & 同上 \\
$\mathrm{LDL}^\top$ 求解 & {\bfseries FP64} & 除法对 $D_j$ 高度敏感 \\
位姿误差/阻尼 & FP64 & 精度关键路径 \\
\bottomrule
\end{tabular}
\end{table*}

\subsubsection{Jacobian 精度验证与混合精度边界}

对 500 个随机 UR10 关节构型，分别以 FP64 和 FP32 计算中心差分 Jacobian，比较相对 Frobenius 误差。

\begin{table*}[ht]
\centering\small
\caption{FP32 Jacobian 精度验证（500 随机关节构型）}
\label{tab:jacobian}
\begin{tabular}{ccccc}
\toprule
$\varepsilon$ & 中位误差 & 均值误差 & p95 误差 & 最大误差 \\
\midrule
$10^{-4}$ & $8.1\times10^{-4}$ & $7.9\times10^{-4}$ & $1.0\times10^{-3}$ & $1.2\times10^{-3}$ \\
$10^{-5}$ & $1.9\times10^{-3}$ & $2.0\times10^{-3}$ & $2.7\times10^{-3}$ & $3.6\times10^{-3}$ \\
$10^{-6}$ & $4.2\times10^{-2}$ & $4.2\times10^{-2}$ & $5.1\times10^{-2}$ & $6.4\times10^{-2}$ \\
$10^{-7}$ & $7.0\times10^{-1}$ & $6.8\times10^{-1}$ & $9.1\times10^{-1}$ & $1.0\times10^{0}$ \\
\bottomrule
\end{tabular}
\end{table*}

实验结果揭示了 FP32 混合精度的精度--速度 Pareto 前沿：

（1）$\varepsilon=10^{-6}$ 下 FP32 Jacobian 约 2 位有效数字（中位误差 $4.2\times10^{-2}$），对应 IK 收敛精度边界约 4--5 mm——解释了后文表~\ref{tab:maincomp} 中 Medium（10 mm）成功而 Strict（5 mm）失败的现象。

（2）$\varepsilon=10^{-5}$ 可将 Jacobian 精度提升一个量级（中位 $1.9\times10^{-3}$），代价是差分步长增大 10 倍。$\varepsilon=10^{-7}$ 发生灾难性抵消（中位 0.70），确认 FP32 物理边界。

（3）保护机制有效：$\mathbf{J}^\top\mathbf{W}^2\mathbf{J}$ 内积在 FP64 归约，$\mathrm{LDL}^\top$ 与收敛判定保留 FP64，消融实验证实 Medium 阈值下收敛率 $\ge0.998$，与全 FP64（1.000）无实质差异。

该分析{\bfseries 为混合精度 IK 的精度--速度 Pareto 前沿提供了量化基准}：吞吐优先场景（料箱抓取、轨迹采样）使用 $\varepsilon=10^{-6}$；更高精度需求场景下 $\varepsilon=10^{-5}$ 或全 FP64 路径的验证将作为后续工作。

\subsection{自适应阻尼策略}

固定阻尼 $\lambda\equiv2\times10^{-3}$ 为基线。三阶段自适应策略[3]：（1）迭代 0 距离初始化——$e_{\mathrm{pos}}>0.5$ m$\to\lambda=0.1$，$0.1<e_{\mathrm{pos}}\le0.5$ m$\to$线性插值，否则$\to5\times10^{-4}$；（2）Marquardt 误差驱动——$r=e^{(k)}_{\mathrm{pos}}/e^{(k-1)}_{\mathrm{pos}}$，$r<0.9\to\lambda\times0.7$，$r>1.1\to\lambda\times2.0$；（3）停滞超驰——连续 12 次无改善$\to\lambda\times5.0$。$\lambda\in[10^{-4},0.5]$。

\section{实验设计}\label{sec:exp}

\subsection{统一 Benchmark 与三层计时口径}

所有目标由 seed=42 确定性生成：随机采样合法 $\mathbf{q}_{\mathrm{gt}}$ 经 FK 生成 $\mathbf{T}_{\mathrm{target}}$。统一初始种子 $\mathbf{q}_{\mathrm{seed}}=[0,0,0,0,0,0]$。所有方法使用同一批目标。

\begin{table*}[ht]
\centering\small
\caption{统一 Benchmark 标准配置}
\label{tab:benchmark}
\begin{tabular}{lll}
\toprule
参数 & 配置 & 理由 \\
\midrule
机器人模型 & UR10 (URDF, tool0) & 代表性 6-DOF 工业机械臂 \\
目标位姿 & seed=42, $N=100\to10\,000$ & 可复现确定性生成 \\
初始种子 & zero\_seed（全零关节角） & 运动学最远端 \\
主收敛阈值 & 10 mm / 5$^\circ$（Medium） & 显著区分求解质量 \\
辅助阈值 & Loose 30mm/10$^\circ$, Strict 5mm/1$^\circ$ & 多级压力测试 \\
最大迭代次数 & 160 & 工业保守上限 \\
碰撞检测 & 关闭（所有方法统一） & 单 seed 场景碰撞无关 \\
重复次数 & 30（含 3 次预热） & 排除冷启动延迟 \\
硬件平台 & RTX 4060 Laptop, CUDA 12.6 & Ada Lovelace, sm\_89 \\
\bottomrule
\end{tabular}
\end{table*}

{\bfseries 对比方法。} CUDA-Mixed（A7，无 CUDA Graph）；CUDA-Mixed-Graph（A8，+CUDA Graph）；cuRobo-NoGraph；cuRobo-Graph。cuRobo 配置：num\_seeds=1，seed\_solver\_num\_seeds=1，self\_collision\_check=False。

{\bfseries 三层计时口径。} 这是本实验设计的关键改进：（1）End-to-End——H2D$\to$kernel$\to$D2H$\to$同步完整延迟；（2）GPU Stream——纯设备侧 CUDA event 时间；（3）Nsight Systems——kernel launch 数、API 调用数等机制分析。所有方法 3 次预热后计时。

{\bfseries 评价指标。} raw\_throughput = $N$/GPU\_stream\_time\_sec；valid\_throughput = raw $\times$ success\_rate（同时衡量速度与质量）；误差分布 p50/p95/max；30 次重复的 mean/std/p50/p95。

\subsection{消融与 Jacobian 精度验证}

消融覆盖 A0--A7 共 8 级。Jacobian 精度验证见 §2.4.1。

\section{实验结果与分析}\label{sec:results}

\subsection{主性能对比}

\begin{table*}[ht]
\centering\small
\caption{主性能对比（Medium 10mm/5$^\circ$，30 次重复，GPU stream 时间）}
\label{tab:maincomp}
\begin{tabular}{clccccc}
\toprule
$N$ & 方法 & Graph & GPU流均值$\pm$std/ms & raw t/s & 成功率 & valid t/s \\
\midrule
100 & CUDA-Mixed & -- & $1.24\pm0.04$ & 80 757 & 1.000 & 80 757 \\
100 & cuRobo & 开 & $2.60\pm0.13$ & 38 398 & 1.000 & 38 398 \\
\midrule
500 & CUDA-Mixed & -- & $3.89\pm0.05$ & 128 404 & 0.998 & 128 147 \\
500 & cuRobo & 开 & $2.60\pm0.14$ & 192 400 & 1.000 & 192 400 \\
\midrule
1 000 & CUDA-Mixed & -- & $7.21\pm0.49$ & 138 719 & 1.000 & 138 719 \\
1 000 & cuRobo & 开 & $2.99\pm0.24$ & 334 189 & 1.000 & 334 189 \\
\midrule
4 000 & CUDA-Mixed & -- & $25.44\pm1.11$ & 157 236 & 1.000 & 157 236 \\
4 000 & cuRobo & 开 & $27.50\pm0.84$ & 145 437 & 0.9998 & 145 401 \\
\midrule
5 000 & CUDA-Mixed & -- & $31.75\pm0.50$ & 157 481 & 0.999 & 157 387 \\
5 000 & cuRobo & 开 & $4.73\pm0.29$ & 1 056 171 & 1.000 & 1 056 171 \\
\midrule
10 000 & CUDA-Mixed & -- & $63.51\pm0.60$ & 157 454 & 0.9995 & 157 375 \\
10 000 & cuRobo & 开 & $61.29\pm1.26$ & 163 160 & 0.9999 & 163 144 \\
\bottomrule
\multicolumn{6}{l}{\footnotesize 注：cuRobo-NoGraph 数据见表~\ref{tab:nsys}。N=5 000 cuRobo-Graph 为 batch-size 甜点，性能不可推广至邻近 N。} \\
\end{tabular}
\end{table*}

表~\ref{tab:maincomp} 和图~\ref{fig:throughput}--\ref{fig:gpu_stream} 展示了主性能对比。

\begin{figure}[ht]
\centering
\includegraphics[width=\columnwidth]{fig2_throughput.png}
\caption{批量吞吐对比（Raw 与 Valid Throughput）}
\label{fig:throughput}
\end{figure}

{\bfseries （1）CUDA Graph 对 cuRobo 影响巨大。} 开启 Graph 后 cuRobo 吞吐提升 7--13 倍。\texttt{use\_cuda\_graph=False} 严重低估其性能。本工作以 cuRobo-Graph 为公平基线，cuRobo-NoGraph 仅作为参考。

{\bfseries （2）$N=100$ 极小批量优势。} N=100 时 CUDA-Mixed GPU 流时间（$1.24\pm0.04$ ms）仅为 cuRobo-Graph（$2.60\pm0.13$ ms）的 47.7\%，延迟降低约 52\%。优势源于单 kernel 零调度开销。但当 $N\ge500$ 时，cuRobo-Graph 的吞吐（192k--334k t/s）显著超过 CUDA-Mixed（128k--157k t/s），因为计算量增大后调度开销占比下降，cuRobo 的粒子优化+L-BFGS 算法效率体现出来。

{\bfseries （3）$N=5\,000$ 是 cuRobo-Graph 的特殊甜点。} 经独立复核实验确认（见实验复核报告），$N=5\,000$ 时 cuRobo-Graph 的 $4.73\pm0.29$ ms 和 1 056k t/s 是真实结果——该 batch size 恰好使 cuRobo 内部批处理+Graph 流水线达到最优。但该性能不可推广：$N=4\,000$ 时 cuRobo-Graph 退化至 $27.50$ ms（5.8 倍慢），$N=10\,000$ 时 $61.29$ ms。CUDA-Mixed 在全部 $N$ 值呈严格线性（$R^2>0.999$），性能可预测。

\begin{figure}[ht]
\centering
\includegraphics[width=\columnwidth]{fig3_e2e.png}
\caption{端到端延迟对比（含 H2D/D2H 传输）}
\label{fig:e2e}
\end{figure}

\begin{figure}[ht]
\centering
\includegraphics[width=\columnwidth]{fig4_gpu_stream.png}
\caption{GPU Stream 执行时间对比（纯设备侧）}
\label{fig:gpu_stream}
\end{figure}

\subsection{求解质量与误差分布}

\begin{table*}[ht]
\centering\small
\caption{位置与姿态误差分布（$N=500$，Medium 阈值）}
\label{tab:errors}
\begin{tabular}{lcccccc}
\toprule
方法 & pos mean/mm & pos p50/mm & pos p95/mm & rot mean/($^\circ$) & rot p50/($^\circ$) & rot p95/($^\circ$) \\
\midrule
CUDA-Mixed & 4.34 & 3.81 & 9.19 & 0.177 & 0.121 & 0.541 \\
cuRobo-Graph & $<0.001$ & $<0.001$ & $<0.001$ & $<0.001$ & $<0.001$ & $<0.001$ \\
\bottomrule
\end{tabular}
\end{table*}

表~\ref{tab:errors} 和图~\ref{fig:errors} 比较了误差分布。cuRobo-Graph 精度达亚微米量级——粒子优化固有优势。CUDA-Mixed 的 p95 位置误差 9.19 mm 在 Medium 阈值内，margin 0.81 mm，中位数 3.81 mm。两方法成功率均 $\ge0.998$。CUDA-Mixed 的意义在于中等精度批量求解，并非高精度最终 IK 求解器的完全替代。

\begin{figure}[ht]
\centering
\includegraphics[width=\columnwidth]{fig5_errors.png}
\caption{误差分布对比（p50/p95 位置与姿态误差）}
\label{fig:errors}
\end{figure}

\subsection{CUDA Graph 对单 Kernel 方案的影响}

\begin{table*}[ht]
\centering\small
\caption{A7 Direct vs A8 CUDA Graph 对比（均值$\pm$标准差）}
\label{tab:graphcmp}
\begin{tabular}{ccccc}
\toprule
$N$ & A7 Direct/(t/s) & A8 Graph/(t/s) & 差异/\% & 判定 \\
\midrule
100 & $80\,757$ & $81\,240$ & +0.6 & 可忽略 \\
500 & $128\,404$ & $142\,805$ & +11.2* & 受 GPU 状态影响 \\
1 000 & $138\,719$ & $147\,381$ & +6.2* & 受 GPU 状态影响 \\
5 000 & $157\,481$ & $157\,841$ & +0.2 & 可忽略 \\
\bottomrule
\multicolumn{5}{l}{\footnotesize *N=500/1\,000 的 A8 标准差偏高（复核确认 GPU 热状态波动），首次实验与复核实验的 A8 相对 A7 提升在 0--11\% 间。} \\
\multicolumn{5}{l}{\footnotesize N$\ge$1\,000 时差异稳定 $<1\%$；N$\le$500 时 CUDA Graph 可能有小到中等幅度影响（0--11\%），} \\
\multicolumn{5}{l}{\footnotesize 但该影响显著小于 cuRobo 从 NoGraph 到 Graph 的 7--13$\times$ 量级提升。} \\
\end{tabular}
\end{table*}

表~\ref{tab:graphcmp} 显示：单 kernel 方案从 CUDA Graph 获益有限，尤其在 $N\ge1\,000$ 时差异 $<1\%$。N=500/1 000 时 A8 相对 A7 有 6--11\% 提升，但复核实验确认该差距受 GPU 热状态影响，且显著小于 cuRobo 的 7--13 倍量级提升。整体上，单 kernel 结构降低了对 CUDA Graph 的依赖。

\subsection{消融实验}

\begin{table*}[ht]
\centering\small
\caption{关键消融结果（Medium 10mm/5$^\circ$）}
\label{tab:ablationdata}
\begin{tabular}{ccccccc}
\toprule
$N$ & A0 基线/(t/s) & 收敛率 & A5 阻尼/(t/s) & 收敛率 & A7 混合/(t/s) & 收敛率 \\
\midrule
100 & 7 589 & 0.830 & 51 361 & 1.000 & 80 757 & 1.000 \\
500 & 9 896 & 0.522 & 62 384 & 0.998 & 128 404 & 0.998 \\
1 000 & -- & -- & -- & -- & 138 719 & 1.000 \\
5 000 & 12 507 & 0.564 & 66 050 & 0.9998 & 157 481 & 0.999 \\
\bottomrule
\end{tabular}
\end{table*}

表~\ref{tab:ablationdata} 和图~\ref{fig:ablation} 展示了逐级优化。自适应阻尼（A0$\to$A5）是收敛率决定性因素：固定阻尼收敛率 52--83\%，A5 恢复至 0.998+。吞吐提升源于迭代次数下降（$N=500$：79.7$\to$13.8，$-83\%$）。混合精度（A5$\to$A7）额外贡献 120--149\% 吞吐，收敛率无退化。A1--A4 独立贡献合计 $<5\%$。

\begin{figure}[ht]
\centering
\includegraphics[width=\columnwidth]{fig6_ablation.png}
\caption{逐级消融吞吐提升（A0$\to$A7）}
\label{fig:ablation}
\end{figure}

\subsection{Nsight Systems 机制分析}

\begin{table*}[ht]
\centering\small
\caption{Nsight Systems CUDA API trace 分析}
\label{tab:nsys}
\begin{tabular}{lcccc}
\toprule
$N$ & 方法 & kernel launches & API calls & GPU active/ms \\
\midrule
100 & CUDA-Mixed & 1 & 5 & 1.24 \\
4 000 & CUDA-Mixed & 1 & 5 & 25.44 \\
4 000 & cuRobo-NoGraph & 14 108 & 28 540 & 197.32 \\
5 000 & cuRobo-NoGraph & 5 945 & 12 014 & 28.16 \\
5 000 & cuRobo-Graph & 2 & 12 & 4.73 \\
\bottomrule
\end{tabular}
\end{table*}

表~\ref{tab:nsys} 和图~\ref{fig:degradation} 解释了 cuRobo 批量振荡根因。$N=4\,000$（退化点）产生 14 108 次 kernel launch——是 $N=5\,000$（正常点）的 2.37 倍。每次 launch 开销 5--10 $\mu$s 累积至 70--140 ms。cuRobo-Graph 将 launch 数降至 2。

cuRobo 退化呈非单调跳变，\texttt{cudaMalloc/cudaFree} 占比 $<1\%$，根因在子批次划分策略的批量依赖性。本文方案对此免疫：launch 数恒为 1。

\begin{figure}[ht]
\centering
\includegraphics[width=\columnwidth]{fig7_degradation.png}
\caption{cuRobo 批量规模退化点分析}
\label{fig:degradation}
\end{figure}

\subsection{Nsight Compute 剖析与批量扩展性}

\begin{table}[ht]
\centering\small
\caption{Nsight Compute 剖析对比（$N=100$）}
\label{tab:ncu}
\begin{tabular}{lccc}
\toprule
指标 & FP64全优化 & CUDA-Mixed & 变化 \\
\midrule
计算吞吐率/\% & 66.89 & 60.73 & $-9$ \\
显存吞吐率/\% & 1.56 & 1.16 & -- \\
寄存器/线程 & 94 & 98 & $+4$ \\
Bank 冲突/次 & 3 522 & 1 295 & {\bfseries $-63\%$} \\
局部内存溢出 & 0 & 0 & -- \\
Kernel 时间/$\mu$s & 2 920 & 827 & {\bfseries $-72\%$} \\
\bottomrule
\end{tabular}
\end{table}

表~\ref{tab:ncu} 和图~\ref{fig:ncu} 确认 kernel 为计算密集型（计算 60--67\% vs 显存 1--2\%），零寄存器溢出。

\begin{figure}[ht]
\centering
\includegraphics[width=\columnwidth]{fig8_ncu.png}
\caption{Nsight Compute 剖析指标对比}
\label{fig:ncu}
\end{figure}

{\bfseries 批量扩展性。} CUDA-Mixed 在 $N=100\to10\,000$ 全量程 GPU 时间严格线性（$R^2>0.999$），$N\ge500$ 时吞吐 128k--157k t/s。$N=100$ 因 block 不足约 81k t/s。批量无关性意味着用户无需对不同 $N$ 值预扫描或参数调优——这是 1 block/target 结构确定性的直接结果。

\section{结 论}

在本文设定的 UR10 单 seed、关闭碰撞检测的批量 IK 查询任务中（统一目标、种子和收敛阈值），自定义 CUDA 小矩阵 DLS 求解器通过 1 block/target 映射、单 kernel 全迭代封装、寄存器级 $\mathrm{LDL}^\top$ 求解和 FP32/FP64 混合精度策略，实现了稳定的线性批量扩展。

实验表明：（1）$N=100$ 极小批量下 GPU 流延迟较 cuRobo-Graph 降低约 52\%；（2）当 $N\ge500$ 时 cuRobo-Graph 依靠 Graph 和内部批处理通常获得更高吞吐，但存在批量依赖的非单调退化（$N=4\,000$ 时退化 5.8 倍）；（3）本文方法 kernel launch 数恒为 1，全量程 $N=100\to10\,000$ GPU 时间严格线性（$R^2>0.999$），性能可预测——这是相对 cuRobo 的核心结构性优势；（4）FP32 混合精度在 $\varepsilon=10^{-6}$ 下存在约 4--5 mm IK 收敛精度边界，Medium 阈值下成功率 0.998+，适用于中等精度批量采样与候选解筛选。Strict 阈值下 $\varepsilon=10^{-5}$ 或全 FP64 路径的精度恢复验证将作为后续工作。

后续工作包括：7 自由度冗余机械臂验证；多初始种子 GPU 并行化；Strict 阈值下 $\varepsilon=10^{-5}$ 配置的精度恢复实验；带碰撞约束的运动规划场景扩展。

\newpage
\begin{thebibliography}{99}
\bibitem{1} SICILIANO B, SCIAVICCO L, VILLANI L, et al. Robotics: modelling, planning and control [M]. 2nd ed. London: Springer, 2010.
\bibitem{2} BUSS S R. Introduction to inverse kinematics with Jacobian transpose, pseudoinverse and damped least squares methods [J]. IEEE Journal of Robotics and Automation, 2004, 17(1): 1--19.
\bibitem{3} MARQUARDT D W. An algorithm for least-squares estimation of nonlinear parameters [J]. Journal of the Society for Industrial and Applied Mathematics, 1963, 11(2): 431--441.
\bibitem{4} 贾龙飞, 乔尚岭, 陶云飞, 等. 冗余机械臂逆运动学求解方法研究进展 [J]. 控制与决策, 2023, 38(12): 3267--3284.\\
JIA L F, QIAO S L, TAO Y F, et al. Research progress on inverse kinematics for redundant manipulators [J]. Control and Decision, 2023, 38(12): 3267--3284.
\bibitem{5} BRUYNINCKX H. Open robot control software: the OROCOS project [C]//Proc. of the IEEE International Conference on Robotics and Automation. Seoul: IEEE, 2001: 2523--2528.
\bibitem{6} SUNDARALINGAM B, HARIH S K S, FISHMAN A, et al. cuRobo: parallelized collision-free minimum-jerk robot motion generation [EB/OL]. [2026-06-15]. https://arxiv.org/abs/2310.17274.
\bibitem{7} YASUTAKE S, KINGSTON Z, PLANCHER B. HJCD-IK: GPU-accelerated inverse kinematics through batched hybrid Jacobian coordinate descent [EB/OL]. [2026-06-15]. https://arxiv.org/abs/2510.07514.
\bibitem{8} ABOELNASR M, et al. ManipulaPy: a GPU-accelerated Python framework for robotic manipulation [J]. Journal of Open Source Software, 2025, 10: 8490.
\bibitem{9} PERRAULT N, HO Q H, LAHIJANIAN M. Kino-PAX: highly parallel kinodynamic sampling-based planner [EB/OL]. [2026-06-15]. https://arxiv.org/abs/2409.06807.
\bibitem{10} ABDELFATTAH A, HAIDAR A, TOMOV S, et al. Batched one-sided factorizations of tiny matrices using GPUs: challenges and countermeasures [J]. Journal of Computational Science, 2018, 26: 226--236.
\bibitem{11} 刘世芳, 赵永华, 黄荣锋, 等. 基于批量LU分解的矩阵求逆在GPU上的有效实现 [J]. 软件学报, 2023, 34(11): 4952--4972.\\
LIU S F, ZHAO Y H, HUANG R F, et al. Efficient GPU implementation of matrix inversion based on batched LU decomposition [J]. Journal of Software, 2023, 34(11): 4952--4972.
\bibitem{12} SHRIDHAR S A. Optimized block-level matrix inversion kernels for small, batched matrices on GPUs [D]. 2024.
\bibitem{13} ABDELFATTAH A, HAIDAR A, TOMOV S, et al. Mixed-precision iterative refinement using tensor cores on GPUs to accelerate solution of linear systems [J]. Proceedings of the Royal Society A, 2021, 477(2253): 20200110.
\bibitem{14} Universal Robots. Universal Robots ROS2 description package [EB/OL]. [2026-06-15]. https://github.com/UniversalRobots/Universal\_Robots\_ROS2\_Description.
\bibitem{15} NVIDIA Corporation. CUDA C++ programming guide [EB/OL]. [2026-06-15]. https://docs.nvidia.com/cuda/cuda-c-programming-guide/.
\end{thebibliography}

\vspace{5mm}
\noindent{\bfseries 作者简介}

\vspace{2mm}
\noindent （匿名稿——作者信息在单独标题页）\\

\end{document}
